\input{../header}

\begin{document}
%


%\onehalfspacing
\allowdisplaybreaks
%##################################################################
\section{PS\#3:  Integration by substitution and integration by parts - \red{Answer key} }

\everymath{\displaystyle}
\begin{enumerate}[leftmargin=0pt]

\item There are some integrals you can do in multiple different ways. Here's one for you to think about.
\begin{enumerate}
    \item Compute $\int\cos(x)\sin(x)\, dx$ using substitution.
    
    \begin{red}
    I think I'm going to say $u = \sin(x)$. Then $\dfrac{du}{dx} = \cos(x)$, so $du = \cos(x)\, dx$. Well how about that, I have $\cos(x)\, dx$ in the integrand!
    
    \[\int\cos(x)\sin(x)\, dx = \int\underbrace{\sin(x)}_{u}\, \underbrace{\cos(x)\, dx}_{du} = \int u\ du = \frac12 u^2 + C = \frac12 (\sin(x))^2 + C\]
    
    (Notice that we could also have used $u = \cos(x)$, $du = -\sin(x)\, dx$. Do we get a different answer that way, or is it really the same answer in disguise?)
    \end{red}
    \item Compute $\int\cos(x)\sin(x)\, dx$ using integration by parts.
    
    \begin{red}
    I think I'll use this combination of $u$ and $dv$:
    \[\begin{array}{ll}
    u = \sin(x) & \\
     & dv = \cos(x)\ dx
    \end{array}
    \textrm{ so then }
    \begin{array}{ll}
        u = \sin(x) & v = \sin(x) \\
        du = \cos(x)\ dx & dv = \cos(x)\ dx
    \end{array}
    \]
    Then using the integration by parts formula, 
    \begin{align*}
        \int\cos(x)\sin(x)\ dx &= \sin(x)\sin(x) - \int\sin(x)\cos(x)\ dx \\
        \intertext{Oh look, the second integral is the same as the first! Let's add it back to the other side.}
        2\int\cos(x)\sin(x)\ dx &= \sin^2(x) \\
        \int\cos(x)\sin(x)\ dx &= \frac12 \sin^2(x) + C
    \end{align*}
    (Notice that this probably would have worked if we did $u = \cos(x)$ and $dv = \sin(x)\ dx$. Do we get a different answer that way, or is it really the same answer in disguise?)
    \end{red}
\end{enumerate}

\item (\href{https://activecalculus.org/single/sec-5-4-parts.html#CBY}{Section 5.4 Exercise 7})
For each of the following indefinite integrals, determine whether you would use 
$u$-substitution, integration by parts, neither, or both to evaluate the integral. In each case, write one sentence to explain your reasoning, and include a statement of any substitutions used. Finally, use your chosen approach to evaluate each integral.
\begin{enumerate}
    \item $\int x^2 \cos(x^3)\ dx$
    
    \begin{red}
    This one is pretty clearly substitution, because the $x^2$ looks like the derivative of the $x^3$ filling. Indeed, if $u = x^3$, then $\frac{du}{dx} = 3x^2$, so $dx = \frac{du}{3x^2}$.
    
    \begin{align*}
        \int x^2 \cos(x^3)\ dx &= \int x^2 \cos(u)\ \frac{du}{3x^2} = \frac13 \int \cos(u)\ du = \frac13 \sin(u) + C = \frac13 \sin(x^3) +C
    \end{align*}
    \end{red}
    
    \item $\int x^5 \cos(x^3)\ dx = \int x^3\cdot (x^2\cos(x^3))\ dx$
    
    \begin{red}
    This one we'll use integration by parts, specifically because we just figured out how to integrate $x^2\cos(x^3)$. 
    \[
    \begin{array}{ll}
        u = x^3 &  \\
         & dv = x^2\cos(x^3)\ dx
    \end{array}
    \textrm{ so then }
    \begin{array}{ll}
        u = x^3 &  v = \tfrac{1}{3} \sin(x^3)\\
        du = 3x^2\ dx & dv = x^2 \cos(x^3)\ dx 
    \end{array}
    \]
    Using the integration by parts formula:
    \begin{align*}
        \int x^3\cdot (x^2\cos(x^3))\ dx &= x^3\cdot \tfrac{1}{3} \sin(x^3) - \int \tfrac{1}{3} \sin(x^3)\cdot 3x^2\ dx \\
        &= \frac13 x^3 \sin(x^3) - \int\sin(x^3)\cdot x^2\ dx
        \intertext{We'll use the substitution $t = x^3$, $dt = 3x^2\ dx$ to evaluate this second integral.}
        &= \frac13 x^3 \sin(x^3) - \int\sin(t)\ \frac{dt}{3} \\
        &= \frac13 x^3 \sin(x^3) + \frac13 \cos(t) = \frac13 x^3 \sin(x^3) + \frac13 \cos(x^3) + C
    \end{align*}
    \end{red}
    
    \item $\int x\ln(x^2)\ dx$
    
    \begin{red}
    This one looks like substitution to me, because the $x$ on the outside looks like the derivative of the $x^2$. Let's say $u = x^2$, so $\frac{du}{dx} = 2x$, so $dx = \frac{du}{2x}$.
    
    \begin{align*}
        \int x\ln(x^2)\ dx &= \int x \ln(u)\ \frac{du}{2x} = \frac12 \int\ln(u)\ du - \textrm{We figured out that integral in class!}\\
        &= \frac12\left(u\ln(u) - u\right) = \frac12\left(x^2 \ln(x^2) - x^2\right) + C
    \end{align*}
    \end{red}
    
    \begin{blue}
    An alternative strategy using integration by parts: logarithmic functions are often a really nice choice for $u$ because they get so much simpler when you differentiate them.
    
    \[\begin{array}{ll}
        u = \ln(x^2) &  \\
         & dv = x\ dx
    \end{array}
    \textrm{ so }
    \begin{array}{ll}
        u = \ln(x^2) & v = \tfrac12 x^2 \\
        du = \frac{1}{x^2}\cdot 2x\ dx & dv = x\ dx 
    \end{array}
    \]
    
    Applying the integration by parts formula:
    \begin{align*}
        \int x\ln(x^2)\ dx &= \ln(x^2) \cdot \tfrac12 x^2 - 
        \int \frac12 x^2 \cdot \frac{1}{x^2} \cdot 2x\ dx \\
        &= \ln(x^2) \cdot \tfrac12 x^2 - \int x\ dx \textrm{ (!!)} \\
        &= \frac12 x^2 \ln(x^2) - \frac12 x^2 + C
    \end{align*}
    \end{blue}
    
    \item $\int \sin(x^4)\ dx$
    
    \begin{red}
    This function has no elementary antiderivative! $u$-substitution won't work because you can't make up your $du$; parts won't work either.
    \end{red}
    
    \item $\int x^3\sin(x^4)\ dx$
    
    \begin{red}
    Ah ha, now we can use substitution. If $u = x^4$, then $\frac{du}{dx} = 4x^3$, so $dx = \frac{du}{4x^3}$.
    
    \begin{align*}
        \int x^3\sin(x^4)\ dx &= \int x^3 \sin(u)\ \frac{du}{4x^3} = \frac14 \int \sin(u)\ du = \frac14(\cos(u)) = -\frac14\cos(x^4)+ C
    \end{align*}
    \end{red}
    
    \item $\int x^7 \sin(x^4)\ dx = \int x^4\cdot (x^3\sin(x^4))\ dx$
    
    \begin{red}
    We'll use a similar trick as we did in part (b):
    \[\begin{array}{ll}
        u = x^4 &  \\
         & dv = x^3\sin(x^4)\ dx
    \end{array}
    \textrm{ so then }
    \begin{array}{ll}
        u = x^4 & \ \ v = -\tfrac14 \cos(x^4) \\
        du = 4x^3\ dx &  dv = x^3\sin(x^4)\ dx
    \end{array}
    \]
    Applying integration by parts:
    \begin{align*}
        \int x^4\cdot (x^3\sin(x^4))\ dx &= x^4 \cdot \left(-\tfrac14 \cos(x^4)\right) - \int -\tfrac14\cos(x^4)\cdot 4x^3\ dx \\
        &= -\tfrac14 x^4\cos(x^4) + \int x^3\cos(x^4)\ dx\\
        \intertext{We'll use the substitution $t = x^4$, $dt = 4x^3\ dx$ to evaluate the second integral.}
        &= -\frac14 x^4\cos(x^4) + \frac14 \int \cos(t)\ dt
         = -\frac14 x^4\cos(x^4) + \frac14 \sin(t) 
         = -\frac14 x^4\cos(x^4) + \frac14 \sin(x^4) + C
    \end{align*}
    \end{red}
\end{enumerate}


\item You're working on some problems with your friend Leonhard, who writes down the following:
\[\int (4t+3)\sin(2t)\ dt
\]
\[\begin{array}{ll}
    u = \sin(2t) & v = 2t^2 + 3t \\
    du = 2\cos(2t)\ dt & dv = (4t + 3)\ dt 
\end{array}
\]
\[\int (4t+3)\sin(2t)\ dt = \sin(2t)(2t^2+3t) - \int\left[(2t^2+3t)\cdot 2\cos(2t)\right]\ dt
\]
Leonhard is now stuck: he's not sure how to attack this new integral. Can you help him get unstuck and find this integral correctly?

\begin{red}
Leonhard's issue is that he made a bad choice for $u$ and $dv$. In particular, the thing he chose for $dv$ got more complicated (went up by a power of $t$) when he integrated it. He can solve his issue by switching his choice:
\[\begin{array}{ll}
    u = 4t+3 & v = -\tfrac12\cos(2t) \\
    du = 4\ dt & dv = \sin(2t)\ dt 
\end{array}
\]
Then the integration by parts formula works out much nicer (look how much simpler $du$ is!):
\begin{align*}
    \int (4t+3)\sin(2t)\ dt &= (4t+3)\cdot\tfrac{-1}{2}\cos(2t) - \int \tfrac{-1}{2}\cos(2t) \cdot 4\ dt 
    =-(4t+3)\cdot\tfrac12\cos(2t) + 2\int\cos(2t)\ dt \\
    & = -(4t+3)\cdot\tfrac12\cos(2t) + \sin(2t) + C
\end{align*}
\end{red}
\end{enumerate}
\end{document}